\documentclass[10pt]{article}
\usepackage[T1]{fontenc}
\usepackage[utf8]{inputenc}
\usepackage[french]{babel}
\usepackage{amsmath,amssymb,mathrsfs,mathpazo}
\usepackage{geometry}
\geometry{a4paper,margin=2cm}
\usepackage{graphicx,array,tabularx,tabularray,colortbl,multirow,multicol,booktabs,enumitem,pifont,ulem,fancybox,diagbox,cancel}
\usepackage{pstricks,pstricks-add,pst-plot,pst-tree,pst-node,pst-text,pst-func,pst-math,pst-bspline,pst-3dplot,pst-eucl,pst-solides3d}
\let\pscalculate\undefined
\usepackage{pst-bezier}
\usepackage{tikz,pgfplots,tkz-tab}
\usetikzlibrary{arrows,arrows.meta,patterns,calc,positioning,trees,decorations.markings,shapes.geometric,shapes.misc,shapes.arrows,intersections,decorations.pathmorphing,decorations.pathreplacing,backgrounds,shadows,babel}
\pgfplotsset{compat=1.18}
\usepgfplotslibrary{fillbetween}
\usepackage[np]{numprint}
\usepackage{listings,makecell,pas-tableur,twoopt}
\let\barre\undefined
\usepackage[pstricks]{tablvar}
\definecolor{darkgreen}{RGB}{0,100,0}
\providecommand{\red}{\color{red}}
\providecommand{\blue}{\color{blue}}
\providecommand{\green}{\color{green}}
\providecommand{\black}{\color{black}}
\providecommand{\white}{\color{white}}
\providecommand{\gray}{\color{gray}}
\providecommand{\cyan}{\color{cyan}}
\providecommand{\magenta}{\color{magenta}}
\providecommand{\yellow}{\color{yellow}}
\providecommand{\violet}{\color{violet}}
\providecommand{\orange}{\color{orange}}
\providecommand{\euro}{\texteuro}
\providecommand{\EUR}{\texteuro}
\providecommand{\ieme}{\textsuperscript{e}}
\providecommand{\degres}{\ensuremath{^\circ}}
\providecommand{\degree}{\ensuremath{^\circ}}
\providecommand{\vv}[1]{\overrightarrow{#1}}
\providecommand{\Coord}[1]{#1}
\providecommand{\decosix}{}
\providecommand{\starredbullet}{\ensuremath{\bullet}}
\pagestyle{empty}
\setlength{\parindent}{0pt}
\usepackage{adjustbox}\DeclareUnicodeCharacter{2236}{:}\DeclareUnicodeCharacter{2190}{\ensuremath{\leftarrow}}
\providecommand{\R}{\mathbb{R}}
\providecommand{\N}{\mathbb{N}}
\providecommand{\Z}{\mathbb{Z}}
\providecommand{\Q}{\mathbb{Q}}
\providecommand{\D}{\mathbb{D}}
\providecommand{\C}{\mathbb{C}}
\providecommand{\e}{{\mathrm{e}}}
\providecommand{\vect}[1]{\overrightarrow{#1}}
\providecommand{\vectt}[1]{\overrightarrow{\mathrm{#1}}}
\providecommand{\vvt}[1]{\overrightarrow{\mathrm{#1}}}
\providecommand{\np}[2][]{{#2~#1}}
\providecommand{\vv}[1]{{\overrightarrow{#1}}}
\renewcommand{\d}{{\mathrm{d}}}
\providecommand{\degres}{{^\circ}}
\providecommand{\degre}{{^\circ}}
\providecommand{\blue}{\color{blue}}
\providecommand{\red}{\color{red}}
\providecommand{\enskip}{\quad}
\providecommand{\llbracket}{\lbrack\!\lbrack}
\providecommand{\rrbracket}{\rbrack\!\rbrack}
\providecommand{\euro}{{\text{€}}}
\providecommand{\no}{{\text{n°}}}
\providecommand{\interval}[2]{{[#1;#2]}}
\providecommand{\intervFO}[2]{{[#1;#2[}}
\providecommand{\dotfill}{\ldots}

\providecommand{\ds}{\displaystyle}
\providecommand{\textding}[1]{#1}
\providecommand{\Oij}{\(\mathrm{O};\vec{\imath},\vec{\jmath}\)}
\providecommand{\Oijk}{\(\mathrm{O};\vec{\imath},\vec{\jmath},\vec{k}\)}
\providecommand{\Ouv}{\(\mathrm{O};\vec{u},\vec{v}\)}
\providecommand{\intervFF}[2]{\ensuremath{[#1;#2]}}
\renewcommand{\euro}{\ifmmode\text{\texteuro}\else\texteuro\fi}
\DeclareUnicodeCharacter{00B0}{\ifmmode^{\circ}\else\textdegree\fi}
\DeclareUnicodeCharacter{20AC}{\ifmmode\text{\texteuro}\else\texteuro\fi}
% Énoncés : Éducation nationale. Transcription et corrigés : APMEP, Denis Vergès et contributeurs. Découpage et classement : CoopMaths, Sébastien Lozano, Éric Elter et contributeurs. Adaptation de la présentation : Maths Explorer.
% CC BY-SA : https://coopmaths.fr/www/about/
% https://coopmaths.fr/alea/static/eam/2026/tex/eam_specialite_2026_06_antilles_1_cor.tex
\begin{document}
\noindent\textit{APMEP / CoopMaths. CC BY-SA. Présentation adaptée par Maths Explorer.}\par\medskip
\begin{enumerate}
	\item \begin{enumerate}
		\item En janvier 2026, il s'est passé un an, donc la valeur de la voiture a perdu 10\,\% de sa valeur initiale.

		Appliquer une baisse de 10\,\%, c'est multiplier par :\quad $q = 1 - \dfrac{10}{100} = 1 - 0,1 = 0,9$.

		$\np{10000} \times 0,9 = \np{9000}.$\quad	La valeur de la voiture en janvier 2026 est de \np[\euro]{9000}.

		\item Pour l'année $n+1$, la valeur de la voiture ($u_{n+1}$) est égale à la valeur de l'année précédente ($u_n$), diminuée de 10\,\% (donc multipliée par $q = 0,9$),

		donc on a bien :\quad $u_{n+1} = 0,9u_n$.
	\end{enumerate}

	\item La relation de récurrence ainsi établie est caractéristique d'une suite géométrique.

	Ainsi, $(u_n)$ est une suite géométrique, de premier terme $u_0 = \np{10000}$ et de raison $q = 0,9$.

	\item On applique la formule qui donne le terme général d'une suite géométrique.

	Pour tout $n$ entier naturel, on a donc :\quad $u_n = u_0\times q^n$

	Soit, ici:\quad $u_n = \np{10000}\times 0,9^n$.

	\item On a :\quad $2030 = 2025 + 5$.

	On cherche donc à calculer $u_5$.

	On a donc :\quad $u_5 = \np{10000}\times 0,9^5 = \np{10000} \times \np{0,59049} = \np{5904,9}$.

	En janvier 2030, la valeur de la voiture est de \np[\euro]{5904,9}.

	\item Ici, la variable \texttt{U} contient la valeur \np{10000}, c'est-à-dire la valeur du terme d'indice 0 de la suite $(u_n)$, la variable \texttt{N} contient l'indice correspondant.

	Tant que la valeur contenue dans \texttt{U} n'est pas strictement inférieure à \texttt{A}, c'est-à-dire tant que cette valeur est supérieure ou égale à \texttt{A}, alors on va remplacer dans \texttt{U} la valeur par la valeur du terme suivant dans la suite $(u_n)$, en appliquant la relation de récurrence, et on va mettre à jour l'indice correspondant.

	On obtient ainsi le programme :

\begin{lstlisting}[style=python,numbers=none,xleftmargin=5cm,xrightmargin=5cm]
def seuil(A) :
	N = 0
	U = 10000
	while U >= A :
		U = 0.9*U
		N = N+1
	return N
\end{lstlisting}

	\item l'appel \texttt{seuil(5000)} renvoie 7. Cela signifie que la valeur de la voiture est strictement inférieure à \np[\euro]{5000} pour la première fois après 7 ans, c'est-à-dire en 2032.

	\item Si la voiture perd plus des trois quarts de sa valeur, cela signifie qu'elle vaut désormais strictement moins d'un quart de sa valeur initiale de \np[\euro]{10000}, soit strictement moins de \np[\euro]{2500}.

	D'après l'appel \texttt{seuil(2500)}, ceci arrivera 14 ans après 2025, c'est-à-dire en 2039.
\end{enumerate}

\needspace{10\baselineskip}

\end{document}
